\documentclass[11pt,reqno]{amsart}

\usepackage[T1]{fontenc}
\usepackage{lmodern}
\usepackage{amsmath,amssymb,mathtools}
\usepackage{microtype}
\usepackage{booktabs,tabularx,longtable,array}
\usepackage{enumitem}
\usepackage{xurl}
\usepackage[hidelinks,unicode]{hyperref}
\urlstyle{same}

\newtheorem{theorem}{Theorem}[section]
\newtheorem{lemma}[theorem]{Lemma}
\newtheorem{corollary}[theorem]{Corollary}
\theoremstyle{definition}
\newtheorem{definition}[theorem]{Definition}
\theoremstyle{remark}
\newtheorem{remark}[theorem]{Remark}
\numberwithin{equation}{section}

\DeclareMathOperator{\RN}{RN}
\DeclareMathOperator{\supp}{supp}
\newcommand{\bits}{B}
\newcommand{\eps}{\varepsilon}
\DeclarePairedDelimiter{\ceil}{\lceil}{\rceil}
\DeclarePairedDelimiter{\floor}{\lfloor}{\rfloor}
\DeclarePairedDelimiter{\abs}{\lvert}{\rvert}
\newcolumntype{Y}{>{\raggedright\arraybackslash}X}
\newcolumntype{P}[1]{>{\raggedright\arraybackslash}p{#1}}
\setlist[enumerate]{itemsep=2pt,topsep=5pt,leftmargin=*}
\setlist[itemize]{itemsep=2pt,topsep=5pt,leftmargin=*}
\emergencystretch=1em
\hypersetup{
  pdftitle={Information bounds and exact indices for ordered floating-point addition},
  pdfauthor={Artem D. Popkov},
  pdfsubject={Exact transfer summaries for sequentially rounded addition},
  pdfkeywords={floating-point arithmetic, roundoff, transfer functions, communication complexity, succinct representation}
}

\title[Ordered floating-point addition]{Information bounds and exact indices\protect\\ for ordered floating-point addition}
\author{Artem D. Popkov}
\date{November 7, 2026}
\subjclass[2020]{Primary 65G50, 68Q11; Secondary 68P30}
\keywords{Floating-point arithmetic, ordered accumulation, transfer function,
  communication complexity, succinct representation}

\begin{document}
\raggedbottom

\begin{abstract}
A block of sequentially rounded additions defines a function of its incoming
accumulator. We study the worst-case number of bits needed to describe this
function before the incoming value is known. At significand precision $p\ge4$,
an explicit family of blocks with $2k$ additions realizes
$\binom{2^{p-2}-2}{k}$ distinct functions, even though each block has real sum
zero and returns zero from an initial zero. This gives optimal-order
$\Theta(mp)$ descriptions for $2\le m\le 2^{\alpha p}$, for each fixed
$0<\alpha<1$, when an operand occupies $\Theta(p)$ bits. For $p\ge6$, on a common query
grid whose rounded trajectories remain within two adjacent positive binades,
we obtain the matching bound $\Theta\bigl(m\log_2(1+2^p/m)\bigr)$ for all
$m\ge2$. Its upper bound follows from a residue potential: adjacent output
gaps start at four and have total deviation at most $4m$. A gcd potential
extends this order on the same grid to any fixed number $K$ of adjacent
positive binades, with constants depending on $K$. A two-threshold normal
form for uniform dyadic rounding compositions gives a different exact
index with at most $1+(K-1)m$ pieces. The same construction covers every
finite input of a binary format with $E$ normal binades, including gradual
underflow and absorbing overflow, using at most $2E+1+2Em$ pieces and
$O(\log(E(m+1)+2))$ abstract operations per query. It identifies signed
zeros numerically and excludes NaNs and nonfinite operands. Its construction
does not require an initial full-state table. A second family
uses permutations of one public multiset of $4k$ distinct operands to realize
$k!$ functions when $2^p\ge64k^2$. The asymptotic lower bounds persist under any fixed
pointwise error probability below one half. Exact arithmetic and native
hardware checks support the constructions. These are bounds on standalone
one-way descriptions.
\end{abstract}

\maketitle

\section{Introduction}\label{sec:introduction}

When rounding is required after each specified addition, summation is an
ordered computation. A block of additions can nevertheless be viewed as a
function of its incoming accumulator. Preparing this function before the
preceding block finishes is one possible route to processing blocks in
parallel. Finite-state transition composition provides such a representation
in principle \cite{ref1,ref2}. We ask how much information a short block of rounded additions can place
in that function.

The encoder knows an ordered operand block $a=(a_1,\ldots,a_m)$ but does not
know the incoming accumulator $x$. It prepares a standalone message from which
a decoder must recover the exact sequentially rounded result at any allowed
$x$. Encoding and decoding time are unrestricted. The resulting lower bounds
therefore apply to tables, compressed data structures, circuits, and any other
representation obeying this access contract. Sending the operands and
replaying the block is an elementary upper bound.

This contract distinguishes a summary from the usual sequential computation.
Once $x$ is known, one accumulator suffices to evaluate the block. With
interaction, the decoder can send $x$ to the block owner and receive the
answer. A description prepared without $x$ must instead preserve the behavior
of the block over the entire query set. We study the information required by
that description, without asserting that ordered addition is inherently
sequential.

A lower bound requires arithmetic attainability.
Counting all functions on a finite set does not establish that rounded
additions realize many of them in a bounded number of operations. Conversely,
monotonicity alone allows maps that need not be generated by a short operand
block. We address both sides of this distinction.

Section~\ref{sec:threshold} constructs an injective encoding of threshold
subsets, using two additions per selected threshold. This gives a
$\Theta(mp)$ description bound for sufficiently short blocks at precision
$p$, when an operand occupies $\Theta(p)$ bits. In
Section~\ref{sec:corridor}, a potential on adjacent output gaps gives the
matching bound
\[
 \Theta\!\left(m\log_2\left(1+\frac{2^p}{m}\right)\right)
\]
on a common query grid under a two-binade trajectory restriction. The same
argument yields a sparse evaluator and a construction algorithm that does
not enumerate the input grid. Section~\ref{sec:multibinade} extends the
potential to every fixed number of adjacent positive binades, using residue
classes of the original grid. Section~\ref{sec:staircase} removes the
exact-sum restriction and gives a different evaluator whose record count
is polynomial in the number of binades. Section~\ref{sec:full-format}
extends that evaluator to the entire finite format, including negative
inputs, subnormals, and absorbing overflow. Section~\ref{sec:order} holds the numerical
multiset fixed and encodes permutations, isolating the information carried
by operand order. Section~\ref{sec:random} extends the lower bounds to
pointwise bounded error by standard entropy arguments.

The matching information bounds retain a trajectory restriction: rounded
states must remain in the specified positive corridor for every allowed
query. We first prove the bounds under a stronger exact-sum promise and
then remove that extra condition. The whole-format evaluator needs no
corridor promise, but has a different, generally larger size bound.
Construction from operands and query evaluation are established; an
efficient merge algorithm on two constructed indices is not claimed.

\section{Related work}\label{sec:related}

\subsection{Transition summaries and ordered accumulation}

Ladner and Fischer compose input-induced maps in a finite-state
construction~\cite{ref1}. Mytkowicz et al. enumerate entry states and exploit
coinciding images and their organization into
ranges~\cite[Sections~3 and~5]{ref2}. Brandoit et al. formulate finite-bit scan
representations and the corresponding cardinality
ceiling~\cite[Section~5 and Appendices~F.4 and~G.2]{ref3}. Our interface
requires only evaluation at the eventual incoming state, and is weaker than
requiring a closed exact merge operation. Transition composition and the
message-counting principle are therefore established tools. The arithmetic
question is how many maps a bounded number of rounded additions can attain.

Kapre and DeHon compute the sequentially specified sum using optimistic
prefixes, recurrence checks, and iterative
correction~\cite[Section~III]{ref4}. Their method retains the original data and
performs further work after detecting errors. It falls outside the standalone
description contract and is not excluded by our bounds. More generally,
nonassociativity alone does not imply a parallel lower bound: saturated
accumulation admits compact associative reformulations~\cite{ref5}.

\subsection{Rounded translations and their compositions}

The rounded translate-and-undo operation is classical. Under stable nearest
rounding, Reiser and Knuth prove that
\[
 T_y(x)=\RN\bigl(\RN(x+y)-y\bigr)
 \quad\text{satisfies}\quad
 T_y(T_y(x))=T_y(x)
\]
whenever the operations are defined~\cite[p.~86, Theorem~1]{ref6}; the
corrections in~\cite{ref7} are incorporated in the later
reprint~\cite{ref18}. Their stable class permits a consistent choice of parity
for each exponent and sign and is broader than a single global ties-to-even
rule. Their second theorem concerns the intervening addition outputs, which
need not already be fixed. Goldberg reproduces the fixed-offset result for
ties-to-even~\cite[Theorem~5]{ref8}. Vor{\'a}{\v c}ek and Hein use related
finite-precision structure in sound randomized smoothing and cite the earlier
result~\cite[Section~3.1]{ref9}.

The local parity-sensitive cancellation in our construction is a
specialization of this classical arithmetic. The exponent-alignment and
midpoint cases of Reiser and Knuth~\cite[pp.~85--86]{ref6} account for its
individual pairs. Our lower bound uses an ordered composition of pairs with
varying offsets and alternating parity to encode a threshold subset.
Section~\ref{sec:order} adds exact surrounding translations and disjoint
threshold bands to encode permutations of one public numerical multiset. On
the witness states, a four-operation slot is itself a rounded translation
$T_{s+r}$; its role is to implement the assignment of thresholds while
preserving the multiset. Neither the primitive nor its idempotence is a new
claim.

The addendum to the reprint points to a subsequent result of
Diamond~\cite[p.~348]{ref18}. For a fixed continuous, strictly monotone function
$f$ that is strictly convex or strictly concave, Diamond proves $h^3=h^2$ on representable inputs for
\[
 h=\RN\circ f^{-1}\circ\RN\circ f,
\]
provided the radix raised to the precision is at least three and all iterates
are defined~\cite[pp.~228--232]{ref19}. Nearest-rounding ties may be chosen
arbitrarily. The proof uses rounding cells and secant slopes along this linked
orbit. Our upper bound instead compares independently initialized trajectories
under varying additions and sums their absolute gap changes over both inputs
and time. Affine translations also lack Diamond's strict curvature
hypothesis. The addendum's multiplication/division pointer concerns another
fixed forward/inverse pair: Knuth's exercise~4.2.2--22 and its answer show
that $M_v(x)=\RN(\RN(xv)/v)$ satisfies $M_v^2=M_v$ on representable inputs
$x$, for fixed nonzero representable $v$, provided no intermediate operation
underflows or overflows~\cite[pp.~229 and~573]{ref23}. The third edition
retains this result with an amended boundary-case
argument~\cite[pp.~244 and~616]{ref20}. Neither version counts the transfer
maps attainable by varying blocks of additions or bounds their description
size.

Casacuberta et al. establish order sensitivity and repeated-rounding effects
in differential-privacy implementations~\cite[Sections~5.3, 5.5, and~5.6]{ref10}.
FPRev infers an unknown accumulation tree by changing the placement of large
positive and negative masks among its inputs~\cite{ref11}; its basic probes
can preserve the multiset. In our interface, the operands and their order in
each block stay fixed while only the incoming scalar varies. Monroe and Job
count structurally inequivalent labeled summation
forms~\cite[Sections~2.1 and~6.1, Proposition~27]{ref17}. That count does not
establish distinct one-variable transfers for one fixed numerical multiset.

Nowak, Cotterell, and Boumasmoud study finite-precision transition monoids and
their aperiodicity~\cite[Appendix~J]{ref12}. Their subsequent transformer
analysis realizes reset behavior on a restricted state action through a
reset-monoid divisor~\cite[Appendix~C]{ref13}. These algebraic properties are
compatible with many distinct length-bounded maps. In particular, the
existence of a rounded reset does not identify the separated query markers
in our attainable families.

\subsection{Compression and information bounds}

Breakpoint propagation under monotone composition is
classical~\cite[Lemma~2.1]{ref14}. Gap coding, subtraction of a known spacing,
and succinct representations of ordered positions are also established
techniques~\cite{ref15}. We prove a linear
bound on the total magnitude of signed gap deviations. Combined with
attainability on the same promised query domain, it gives the matching
description bound. The support estimate and sparse constructor follow from
spacing-four invariance away from a single boundary crossing.

Generic binary rounding theory, including representability and the ordering
of lattice rounding, is formalized in libraries such as FloatLib, building
on Flocq~\cite[Section~2.3]{ref21}. The uniform-lattice facts used below are
elementary ingredients of that theory. Conditions for simplifying consecutive
rounding operations are a different question~\cite[Section~5.2 and
Appendix~A.2]{ref22}: the gap argument retains every prescribed rounding and
compares two trajectories. Its additional statement is a budget for their
total gap variation, charged to crossings of binade boundaries.

The bounded-error extensions use standard augmented-index and entropy
arguments~\cite{ref16}. We include the entropy proof to distinguish that
general inequality from the arithmetic query reductions. The contribution
lies in the attainable families and the aggregate gap bound, rather than in
a new communication inequality or a new rounded cancellation primitive.

\subsection{Staircase representations}

Eggleton classifies fixed insertion patterns of floors and
ceilings in a linear expression~\cite{ref24}; Lagarias et al. classify
commuting dilated floor functions~\cite{ref25}. Our normal form instead
allows arbitrarily long translated nearest-even chains with varying dyadic
spacings. Its proof combines translation equivariance with a common output
coset to leave at most two thresholds in one period.

Monotone piecewise-constant lists and their operations are established
data structures. Henzinger et al. control their size through
approximation~\cite[Section~2]{ref26}. Piecewise symbolic floating-point
models can also deliberately overapproximate individual roundings to
verify programs~\cite[Sections~4.2--4.3]{ref27}. Bagnara et al. treat
IEEE special values and optimal or sound interval projections in
detail~\cite[Sections~1.4 and~4.1--4.3]{ref28}. These contracts differ
from preserving the entire input-output graph of a fixed operand chain.
The index below is exact on every permitted scalar input. Its structural
ingredients are the two-threshold closure and the global count
of preimage cuts, rather than interval lists or IEEE case distinctions
themselves.

\section{Model and notation}\label{sec:model}

All logarithms and information quantities use base two. We write
$\binom{n}{k}$ for a binomial coefficient.

Fix a binary floating-point format with $p$ significant bits, including the
leading bit, and round-to-nearest with ties to even. Let $\RN$ denote correct
rounding. Assume that the finite normal exponent range contains exponents
$0$ through $p$, with the exponent defined by the normalized form
$(1.f)_2\times2^e$. Set $R=2^p$. In the corridor
\[
 C=[R/2,2R),
\]
the lower binade has spacing one and the upper binade has spacing two.
Every representable state in $C$ is an integer. The witness trajectories avoid
overflow, underflow, NaNs, and distinctions between signed zeros. The decoder
returns the numerical result; floating-point exception flags are outside the
contract.

For a block $a=(a_1,\ldots,a_m)$ of finite representable operands, define
\begin{equation}\label{eq:fold}
 s_0=x,\qquad s_i=\RN(s_{i-1}+a_i)\quad(1\le i\le m),\qquad F_a(x)=s_m.
\end{equation}
The format, block length, and allowed query set $X$ are public. An encoder
$E$ depends only on $a$. A decoder $D$ must satisfy
\begin{equation}\label{eq:contract}
 D(E(a),x)=F_a(x)
 \qquad\text{for every allowed block $a$ and every $x\in X$.}
\end{equation}
The message contains all block-specific information available to $D$. A
pointer therefore does not constitute a short standalone summary if it grants
access to uncounted operands or auxiliary data. Public constants and
algorithms independent of the chosen block are free.

Let $B(\mathcal{A},X)$ be the minimum worst-case fixed message length for a
block class $\mathcal{A}$ on $X$. Encoding and decoding are computationally
unrestricted. If the class realizes $M$ distinct functions on $X$, then
\begin{equation}\label{eq:counting}
 B(\mathcal{A},X)\ge \lceil\log_2 M\rceil,
\end{equation}
since different functions cannot share a message. Conversely, an index into
the finite transfer family is an information-theoretic upper bound. These
standard facts depend on the message convention. With arbitrary strings of
length at most $b$ and free external framing, there are $2^{b+1}-1$ messages,
giving the corresponding lower bound $\lceil\log_2(M+1)\rceil-1$.

For bounded-error protocols, public random coins $U$, independent of the
input, may affect both encoding and decoding. For each fixed pair $(a,x)$,
the probability of returning the exact result must be at least $1-\varepsilon$,
where $0\le\varepsilon<1/2$. This is a pointwise guarantee; it does not require
one random message to answer every query correctly at once. Making private
coins public only strengthens a protocol, so our lower bounds also apply to
private randomness.

The corridor class $\mathcal{A}(p,m,X)$ consists of blocks for which both
$s_{i-1}+a_i$ and $s_i$ belong to $C$ at every step and for every $x\in X$.
The encoder is required to describe only blocks satisfying this promise.
The reference constructor also detects violations by checking the extreme
trajectories.

\clearpage
\begin{samepage}
\section{An attainable family of rounded-add transfers}
\label{sec:threshold}

\subsection{Threshold construction}

\begin{theorem}[Binomial attainability]
\label{thm:binomial}
Let $p\geq 4$, put
\[
  X_0=\{8,12,\ldots,R-4\},\qquad N=R/4-2,
\]
and let $0\leq k\leq\min\{N,\lfloor m/2\rfloor\}$ be an integer.
Blocks of length $m$ realize at least $\binom Nk$ distinct transfer
functions on $X_0$. Every block in the constructed family has exact real
sum zero and satisfies $F_a(0)=0$. On $X_0$, all its intermediate values
are finite, normal, and positive. Consequently, every exact fixed-width
description for a class containing this family requires at least
$\lceil\log_2\binom Nk\rceil$ bits.
\end{theorem}
\end{samepage}

\begin{proof}
Choose a $k$-element subset $S\subseteq X_0$ and write its elements as
$t_0<\cdots<t_{k-1}$. Define
\begin{equation}
\label{eq:threshold-radii}
 c_j=
 \begin{cases}
  1,&j\ \text{even},\\
  4,&j\ \text{odd},
 \end{cases}
 \qquad r_j=R+c_j-t_j.
\end{equation}
Concatenate the pairs $(+r_j,-r_j)$ in increasing order of $j$ and
append $m-2k$ positive zeros. Each $r_j$ is an integer strictly between
zero and $R$, so every operand is representable. Every pair has exact
real sum zero and returns exactly to zero when started from zero.

After any prefix of pairs, we show that the state
at a marker $x\in X_0$ is $x-d$, where $d$ is the parity of the number
of processed thresholds not exceeding $x$. Before pair $j$, every marker
$x\geq t_j$ has crossed all $j$ preceding thresholds, and hence
$d=j\bmod 2$. A marker below $t_j$ may have either parity.

If $x<t_j$, the spacing-four grid gives $x\leq t_j-4$. The exact first
sum is the positive integer
\[
  x-d+R+c_j-t_j\leq R.
\]
The subsequent exact difference is $x-d\in[7,R)$. Both operations are
exact, so the state is unchanged.

If $x\geq t_j$ and $j$ is even, then $d=0$. The exact first sum is
$R+1+(x-t_j)$, which lies in $[R,2R)$ and is congruent to one modulo
four. Round-to-nearest with ties to even decreases it by one. Exact
subtraction of $r_j$ then gives $x-1$. If $j$ is odd, then $d=1$; the
exact first sum is $R+3+(x-t_j)$ and is congruent to three modulo four.
It rounds up by one, and subtracting $r_j$ gives $x$. The exact sums
and their rounded values remain strictly below $2R$. Thus the completed
block has transfer function
\begin{equation}
\label{eq:threshold-transfer}
 F_S(x)=x-\bigl(\lvert S\cap(-\infty,x]\rvert\bmod 2\bigr).
\end{equation}

The first threshold in the symmetric difference of two subsets
distinguishes their transfer functions. Equivalently, adjacent
differences of the bits $x-F_S(x)$ recover membership in $S$, with
initial parity zero before the first marker. Hence all $\binom Nk$
transfers are distinct. Padding by positive zero preserves their values.
The fixed-width counting bound completes the proof.
\end{proof}

Adjacent markers differ by four, whereas each final output differs from
its marker by at most one. Their outputs are therefore distinct. Since
deterministic trajectories cannot separate after they have coalesced,
no two marker trajectories in this construction coalesce at an earlier
step. Merging coincident images cannot reduce this family.

\begin{corollary}[Short-block incompressibility]
\label{cor:short}
Fix $0<\alpha<1$. As $p$ tends to infinity, exact summaries require
$\Omega(mp)$ bits, uniformly for $2\leq m\leq 2^{\alpha p}$.
If an operand word has width $w=\Theta(p)$, the optimal worst-case
description length is $\Theta(mp)$.
\end{corollary}

\begin{proof}
For sufficiently large $p$, take $k=\lfloor m/2\rfloor\geq m/3$ and
use $N\geq 2^{p-3}$. The bound $\binom Nk\geq(N/k)^k$ gives
\[
 \log_2\binom Nk\geq k\bigl((1-\alpha)p-3\bigr).
\]
The lower-bound constant may depend on $\alpha$. Sending all $m$
operand words gives the upper bound $mw$.
\end{proof}

For longer blocks, the choice
$k=\min\{\lfloor m/2\rfloor,\lfloor N/2\rfloor\}$ gives a lower
bound of order $m\log_2(1+R/m)$, which saturates at order $R$.
In particular, the bound is not linear in $m$ when the format is fixed
and $m$ tends to infinity. With $m=256$, the exact binomial expression
gives lower bounds of $410$, $2100$, and $5812$ bits at precisions
$p=11$, $24$, and $53$, respectively.

\subsection{The same construction in a narrow corridor}
\label{subsec:threshold-corridor}

For $p\geq 6$, define
\begin{equation}
\label{eq:corridor-grid}
 X_\ast=\{R/2+8,R/2+12,\ldots,3R/4-4\},
 \qquad N_\ast=R/16-2.
\end{equation}
Restrict both thresholds and query markers to $X_\ast$. The preceding
proof still applies. Operand magnitudes lie in $[R/4,R/2)$, every state
at a pair boundary is at least $R/2+7$, and every exact sum and rounded
intermediate value stays in $C=[R/2,2R)$. The family therefore belongs
to $\mathcal{A}(p,m,X_\ast)$, on the same query set for every block.
It has $\binom{N_\ast}{k}$ distinct transfers. Since $N_\ast=\Theta(R)$,
its cardinality has the same asymptotic growth as that of the wider-grid
family.

\section{A matching upper bound in two adjacent binades}
\label{sec:corridor}

\subsection{Local gap dynamics}

Let $x_0,\ldots,x_{L-1}$ be an increasing arithmetic progression of
representable integers with spacing four. Consider a block satisfying
the corridor promise on the entire progression. After a prefix of the
block, write
\[
 y_i=F(x_i),\qquad g_i=y_{i+1}-y_i\quad(0\leq i<L-1).
\]
Correct rounding is monotone, so the gaps are nonnegative integers.
Initially every gap equals four. Adding a common operand preserves
each exact gap before rounding.

\begin{samepage}
\begin{lemma}[Local weighted bound]
\label{lem:weighted}
For a nonnegative integer $g$, define
\begin{equation}
\label{eq:residue-potential}
 \psi(g)=
 \begin{cases}
  0,&g\equiv 0\pmod 4,\\
  2,&g\equiv 2\pmod 4,\\
  3,&g\ \text{odd}.
 \end{cases}
\end{equation}
Suppose one rounded addition changes a gap from $g$ to $h$ and obeys
the corridor promise at both endpoints. If the exact endpoints do not
strictly straddle $R$, then
\begin{equation}
\label{eq:local-charge}
  |h-g|+\psi(h)-\psi(g)\leq 0.
\end{equation}
If they strictly straddle $R$, the left-hand side is at most four.
A gap equal to four is unchanged unless it strictly straddles $R$.
\end{lemma}
\end{samepage}

\begin{proof}
Below $R$, the rounding lattice has spacing one. Two exact values at
an even integer distance round to values at the same distance. For an
odd distance, the gap can change only at a half-integer tie. The new
gap is then even and differs by one; its potential drops from three
to either two or zero, paying for the change.

Above $R$, the lattice has spacing two. Translation by four preserves
ties-to-even rounding, so a gap divisible by four is unchanged. An odd
gap becomes even and changes by one, again paid for by the drop in
potential. A gap congruent to two modulo four can change only at a tie.
When it does, its change has magnitude two and its new value is
divisible by four, so the potential decreases by exactly two.

For a straddling pair, the rounding errors have magnitudes at most
$1/2$ and $1$. Their difference changes the gap by at most $3/2$.
Both the old and new gaps are integers, so $|h-g|\leq 1$. The
potential can increase by at most three, giving a total charge of at
most four. An endpoint equal to $R$ can be included in the
uniform-lattice case containing the other endpoint; the two lattice
descriptions agree at $R$.
\end{proof}

\begin{theorem}[Sparse bounded variation]
\label{thm:variation}
After $m$ additions under the preceding corridor promise, let
$d_i=g_i-4$. Then
\begin{equation}
\label{eq:variation-budget}
 \#\{i:d_i\neq 0\}\leq m,
 \qquad
 \sum_{i=0}^{L-2}|d_i|\leq 4m.
\end{equation}
\end{theorem}

\begin{proof}
At any step, monotonicity permits at most one adjacent pair of exact
endpoints to strictly straddle $R$. Sum the charge from
Lemma~\ref{lem:weighted} over all gaps. The sum of absolute gap changes
plus the change in total potential is at most four. Summing over $m$
steps telescopes. The initial potential is zero and the final potential
is nonnegative, so the total variation over all gap histories is at
most $4m$. The final distance from the initial gap vector is no larger
than this total variation.

For the support bound, an index with zero deviation has gap four. It
can first acquire a nonzero deviation only at a crossing. At most one
new index is introduced per step, proving the other inequality.
\end{proof}

The asymptotic description bound does not require the numerical constant
four to be optimal.

\subsection{Encoding and tightness}

\begin{samepage}
\begin{theorem}[Tight corridor description complexity]
\label{thm:corridor}
For the grid $X_\ast$ in \eqref{eq:corridor-grid}, every $p\geq 6$, and
every $m\geq 2$,
\begin{equation}
\label{eq:tight-corridor}
 B\bigl(\mathcal{A}(p,m,X_\ast),X_\ast\bigr)
 =\Theta\!\left(m\log_2\left(1+\frac{2^p}{m}\right)\right).
\end{equation}
The same order holds for protocols with any fixed pointwise error
$0\leq\varepsilon<1/2$.
\end{theorem}
\end{samepage}

\begin{proof}[Deterministic upper bound]
Set $L=|X_\ast|=N_\ast$. The first output has fewer than $2R$ possible
integer values and can be encoded using $p+1$ bits. It and the
deviations $d_i$ determine all remaining outputs. Put $n=L-1$.
There are $\binom{n+4m}{4m}$ nonnegative integer vectors of length $n$
with sum at most $4m$. At most $\min\{n,4m\}$ nonzero coordinates
require sign bits. A fixed-width index into these descriptions
therefore uses at most
\begin{equation}
\label{eq:enumerative-budget}
 p+1+\left\lceil\log_2\binom{L-1+4m}{4m}\right\rceil
       +\min\{L-1,4m\}
\end{equation}
bits.

For $m\leq R$ and $L=\Theta(R)$, standard binomial bounds give
$O(m\log_2(1+R/m))$ bits; the $p$ term is absorbed when $m\geq 2$.
For $m>R$, monotonicity gives another code. Encode the nondecreasing
integer output ranks by a stars-and-bars bitvector of length
$3R/2+L-1$. This uses $O(R)$ bits, while
$m\log_2(1+R/m)=\Theta(R)$ in this regime. The public parameters
determine which code is used, so this choice requires no block-specific
information.
\end{proof}

\begin{proof}[Deterministic lower bound]
Apply Theorem~\ref{thm:binomial} on $X_\ast$ with
\[
 k=\min\{\lfloor m/2\rfloor,\lfloor N_\ast/2\rfloor\}.
\]
For $m\leq N_\ast$, we have $k\geq m/3$, and
$\binom{N_\ast}{k}\geq(N_\ast/k)^k$ yields
$\Omega(m\log_2(1+N_\ast/m))$ bits. For $m\geq N_\ast$, a central
binomial coefficient has logarithm $\Omega(N_\ast)$. Since
$N_\ast=\Theta(R)$, these regimes give the lower half of
\eqref{eq:tight-corridor}; adjusting universal constants covers the
smallest permitted precisions. Section~\ref{sec:random} proves the
bounded-error lower bound. The deterministic code already supplies
the corresponding upper bound.
\end{proof}

\subsection{Constructing the sparse transfer}
\label{subsec:sparse-construction}

The upper representation can be constructed without evaluating all
$L$ inputs. Maintain the first output $y_0$ and the sorted nonzero
pairs $(j,d_j)$, together with their prefix sums. Then
\begin{equation}
\label{eq:sparse-recovery}
 y_i=y_0+4i+\sum_{j<i}d_j.
\end{equation}
When the support has size $q$, a query takes $O(\log(q+2))$ comparisons
and arithmetic operations. One operand $a$ is processed as follows.
\begin{enumerate}
 \item Evaluate the first and last exact sums and their rounded
       outputs. Reject if either extreme violates the corridor promise.
 \item Find by binary search the first index $i$ for which
       $y_i+a\geq R$, taking $i=L$ if no such index exists.
 \item Merge the sorted deviation indices with the candidate $i-1$
       when $0<i<L$, removing a duplicate candidate if necessary.
 \item For each candidate $j$, compute
       \[
         h=\RN(y_{j+1}+a)-\RN(y_j+a).
       \]
       Retain the new deviation $h-4$ precisely when it is nonzero.
 \item Replace $y_0$ by $\RN(y_0+a)$ and rebuild the prefix sums.
\end{enumerate}

By Lemma~\ref{lem:weighted}, only the candidate gaps can change.
Monotonicity makes the extreme checks sufficient for all exact and
rounded trajectories. The update is therefore exact and detects a
failed promise without enumerating $X$. After $i$ operands its support
has size at most $i$. A deterministic sorted-list implementation uses
\begin{equation}
\label{eq:construction-cost}
 O\!\left(m^2\log(m+2)+m\log L\log(m+2)\right)
\end{equation}
arithmetic, comparison, and rounding operations for construction.
This count excludes serialization and arbitrary-precision bit costs.
A sorted merge suffices, so the bound does not depend on average-case
hash-table behavior.

The reference serializer uses combinatorial subset ranks for deviation
positions and cumulative positive magnitudes, followed by sign bits
and the first value. It pads to the largest payload over all support
sizes at the same public parameters, and uses the monotone bitvector
when that budget is smaller. For compatibility with retained byte
artifacts, the tested serializer uses a conservative $5m$ magnitude
budget while separately checking the stronger $4m$ invariant.
The representation is an evaluator on the specified grid. Closure
under arbitrary composition does not follow: one block's outputs can
leave the next block's query grid or its residue class.

\section{A fixed number of adjacent binades}
\label{sec:multibinade}

The two-binade restriction can be relaxed while retaining the original query
set. Let $K\geq2$ and, for this section, assume a normal exponent range
containing $0$ through $p+K-1$. Define
\begin{equation}
\label{eq:multibinade-parameters}
 C_K=[R/2,2^{K-1}R),\qquad H=2^{K-1},\qquad Q=2H=2^K.
\end{equation}
The spacings in these $K$ binades are $1,2,\ldots,H$. Let
$\mathcal A_K(p,m,X)$ impose the same exact-sum and rounded-state promise
as $\mathcal A(p,m,X)$, with $C_K$ replacing $C$. In particular,
$\mathcal A_2=\mathcal A$. The exponent assumption also makes the upper
boundary representable; valid trajectories must still remain strictly below it.

\subsection{A potential for nested binary lattices}

For every nonnegative integer gap, set
\begin{equation}
\label{eq:gcd-potential}
 \psi_Q(g)=Q-\gcd(Q,g),\qquad \gcd(Q,0)=Q.
\end{equation}
For $Q=4$, this is precisely \eqref{eq:residue-potential}.

\begin{lemma}[Uniform-lattice charge]
\label{lem:uniform-Q}
Suppose two exact endpoints at integer distance $g$ are rounded to the same
uniform lattice of spacing $h$, where $h$ is a power of two and $2h$ divides
$Q$. If their rounded distance is $g'$, then
\[
 |g'-g|+\psi_Q(g')-\psi_Q(g)\leq0.
\]
In particular, a gap divisible by $Q$ is preserved.
\end{lemma}

\begin{proof}
If $2h$ divides $g$, translation by $g$ preserves both the lattice and the
tie parity, so $g'=g$. If $h$ divides $g$ but $2h$ does not, non-ties also
preserve the gap. At a tie, $g'=g\pm h$ is divisible by $2h$. The change
has magnitude $h$ and the potential decreases by at least $h$.

Otherwise put $d=\gcd(Q,g)<h$. The new gap is either
$h\lfloor g/h\rfloor$ or $h\lceil g/h\rceil$. The nonzero remainder of $g$
modulo $h$ is a multiple of $d$, so $|g'-g|\leq h-d$. Since $g'$ is a
multiple of $h$, its greatest common divisor with $Q$ is at least $h$.
The potential decrease is therefore at least $h-d$, as required.
\end{proof}

\begin{lemma}[Charge per crossed boundary]
\label{lem:crossing-Q}
Under the $C_K$ promise, let $c$ be the number of interior binade boundaries
strictly between the two exact endpoints of a rounded addition. Then
\begin{equation}
\label{eq:boundary-Q-charge}
 |g'-g|+\psi_Q(g')-\psi_Q(g)\leq Qc.
\end{equation}
\end{lemma}

\begin{proof}
For $c=0$, Lemma~\ref{lem:uniform-Q} applies. An endpoint exactly on a
boundary is consistent with the lattice on either side containing the other
endpoint: the boundary value is representable with even significand.

Every rounding error has magnitude at most $H/2$. Thus $|g'-g|\leq H=Q/2$,
and the potential can increase by at most $Q-1$. The total charge is less
than $3Q/2$, which proves \eqref{eq:boundary-Q-charge} when $c\geq2$.

Suppose $c=1$. If $Q$ does not divide $g$, then $\gcd(Q,g)\leq Q/2$ and
the potential increase is at most $Q/2-1$. The charge is consequently less
than $Q$. If $Q$ divides $g$, the endpoints lie on adjacent rounding
lattices of spacings $h$ and $2h$. Since $2h\leq H=Q/2$, the power of two
$4h$ divides $Q$. Translation by $g$ therefore preserves ties-to-even on
both lattices. Writing the lower exact endpoint as $u$, we obtain
\[
 g'-g=\RN_{2h}(u)-\RN_h(u),
\]
where the subscripts denote uniform-lattice rounding. These two rounded
values differ by $0$ or $\pm h$. In the nonzero case,
$\gcd(Q,g')=h$, so the charge is exactly $h+Q-h=Q$. Otherwise it is zero.
\end{proof}

\begin{samepage}
\begin{theorem}[Variation across finitely many binades]
\label{thm:multibinade-variation}
Start with a representable arithmetic progression of spacing $Q$ inside
$C_K$. Write $g_i(t)$ for its adjacent output gaps after $t$ additions
satisfying the $C_K$ promise. Then
\begin{align}
 \#\{i:g_i(m)\neq Q\}&\leq (K-1)m,\label{eq:multi-support}\\
 \sum_{t=1}^m\sum_i|g_i(t)-g_i(t-1)|&\leq Q(K-1)m,\label{eq:multi-history}\\
 \sum_i|g_i(m)-Q|&\leq Q(K-1)m.\label{eq:multi-final}
\end{align}
\end{theorem}
\end{samepage}

\begin{proof}
At each step the exact endpoints are nondecreasing. Each interior boundary
is therefore strictly straddled by at most one adjacent pair, so the sum of
their crossing counts is at most $K-1$. Sum
Lemma~\ref{lem:crossing-Q} over all pairs and steps. The initial potential
is zero and the final potential is nonnegative, proving
\eqref{eq:multi-history}. The triangle inequality gives
\eqref{eq:multi-final}. A gap equal to $Q$ cannot first change without a
crossing, by Lemma~\ref{lem:uniform-Q}; at most $K-1$ new deviation indices
can be introduced at a step. This proves \eqref{eq:multi-support}.
Zero gaps and coalesced trajectories satisfy the same argument.
\end{proof}

\subsection{The original query grid and a matching bound}

Changing the grid to spacing $Q$ would change the information problem.
Instead, partition the original $X_\ast$ into its nonempty residue classes
modulo $Q$. There are at most $Q/4=2^{K-2}$ such classes, each an arithmetic
progression of spacing $Q$.

\begin{samepage}
\begin{theorem}[Tight description complexity for fixed $K$]
\label{thm:fixed-K}
For every fixed $K\geq2$, uniformly over $p\geq6$ and $m\geq2$,
\begin{equation}
\label{eq:fixed-K-bound}
 B\bigl(\mathcal A_K(p,m,X_\ast),X_\ast\bigr)
 =\Theta_K\!\left(m\log_2\left(1+\frac{2^p}{m}\right)\right).
\end{equation}
The same order holds for every fixed pointwise error probability
$0\leq\varepsilon<1/2$, with constants allowed to depend on $K$ and
$\varepsilon$.
\end{theorem}
\end{samepage}

\begin{proof}
Let class $j$ contain $L_j$ queries and put $n_j=L_j-1$ and
$V=Q(K-1)m$. Its first output uses at most $p+K-1$ bits. By
Theorem~\ref{thm:multibinade-variation}, a magnitude vector for its gap
deviations has sum at most $V$. The stars-and-bars argument used in
\eqref{eq:enumerative-budget} therefore gives a fixed-width description of
this class using at most
\begin{equation}
\label{eq:multi-code}
 p+K-1+\left\lceil\log_2\binom{n_j+V}{V}\right\rceil
       +\min\{n_j,V\}
\end{equation}
bits. Decode magnitudes first, then read signs in index order, with any
remaining sign positions padded. The formula also applies to a singleton
class, for which $n_j=0$.

For fixed $K$ there are a constant number of classes and $V=O_K(m)$.
When $m\leq R$, summing \eqref{eq:multi-code} and using standard binomial
bounds gives $O_K(m\log_2(1+R/m))$ bits. The first-output terms are absorbed
because $m\log_2(1+R/m)\geq\log_2(1+R)>p$ for $m\geq1$.
For $m>R$, there are exactly $KR/2$ representable states in $C_K$.
Monotone output-rank coding on the original grid gives a fallback of
$O_K(R)$ bits, which has the required order in this regime. The choice of
code depends only on public parameters.

Finally, $\mathcal A_2(p,m,X_\ast)\subseteq\mathcal A_K(p,m,X_\ast)$.
The threshold lower bound on precisely this $X_\ast$, and its bounded-error
version in Section~\ref{sec:random}, therefore apply unchanged. The exact
upper code also supplies a bounded-error upper bound.
\end{proof}

The sparse construction extends as well. Maintain one anchor and a sparse
gap map for each residue class. At each addition, binary-search each of the
$K-1$ boundaries, and update the old deviation indices together with the
resulting crossing candidates. Lemma~\ref{lem:uniform-Q} proves that no
other index can change. For fixed $K$, a sorted-list implementation has the
same asymptotic construction bound as \eqref{eq:construction-cost}, with
constants depending on $K$, and still avoids full-grid enumeration. The
artifact includes this evaluator; it does not include a new serialized
fixed-width codec for arbitrary $K$.

The fixed-$K$ qualification is essential to the stated estimate. This proof
uses up to $2^{K-2}$ residue classes and a variation budget proportional to
$2^K(K-1)m$ per class. It does not give an exponent-range-independent
constant when $K$ grows. Theorem~\ref{thm:fixed-K} is an extension of the
same gap-potential argument, not a claim of unrestricted exact scan closure.

Some dependence on the exponent range is necessary even independently of
this proof. Fix $p=6$, $m=2$, so that $X_\ast=\{40,44\}$. The zero block
and the $K-1$ blocks $(2^{6+j},0)$, $0\leq j\leq K-2$, satisfy the
$C_K$ promise. For each nonzero block put $a=2^{6+j}\geq64$. Its exact
first sums lie in $[a,a+44]$, strictly below $2a-a/64$, the midpoint
between the upper boundary and the preceding representable value.
Thus its rounded outputs remain in the binade $[a,2a)$. These $K$ blocks
have distinct outputs already at the query 40, giving the additional
lower bound $\lceil\log_2 K\rceil$. The expression
$m\log_2(1+2^p/m)$ is constant in this example, so an upper bound with
a constant independent of $K$ cannot hold as $K$ tends to infinity.

\section{Rounded-state corridors and staircase evaluators}
\label{sec:staircase}

The exact-sum restriction can be removed without changing the preceding
information bounds. A different representation also avoids the exponential
number of residue classes in the earlier constructor. Removing the exact-sum restriction
enlarges the promised class; the new representation controls the number
of records needed for exact evaluation when $K$ grows. It does not give a matching information bound uniform in $K$.

\subsection{A strictly weaker trajectory promise}

Let $\mathcal A_K^{\mathrm{rd}}(p,m,X)$ require only that every rounded
state belongs to $C_K$. Retain the format assumptions of
Section~\ref{sec:multibinade}, and put $L=R/2$, $U=2^{K-1}R$ and
$H=2^{K-1}$ in the following lemma.

\begin{lemma}[The permitted rounding cells]\label{lem:rounding-fringe}
For real $z$,
\[
 \RN(z)\in[L,U)
 \quad\Longleftrightarrow\quad L-\tfrac14\leq z<U-H/2.
\]
On $[L-1/4,R]$, floating-point rounding agrees with nearest-even rounding
to the uniform integer lattice.
\end{lemma}

\begin{proof}
The predecessor of $L$ is $L-1/2$, and their midpoint rounds to $L$.
The predecessor of $U$ is $U-H$, and their midpoint rounds to $U$, which
is excluded. Both endpoints have even significands. Monotonicity gives
the asserted preimage. The existing normal exponent range includes these
neighbours. Every point of the additional fringe $[L-1/4,L)$ rounds to
$L$ under both floating-point and uniform integer rounding. Above $L$
and up to $R$, the original unit-lattice description applies.
\end{proof}

\begin{corollary}[Rounded-state description complexity]
\label{cor:rounded-corridor}
The gap-variation and support bounds in
Theorems~\ref{thm:variation} and~\ref{thm:multibinade-variation} hold
unchanged with only the rounded-state promise. For every fixed $K\geq2$,
\[
 B\bigl(\mathcal A_K^{\mathrm{rd}}(p,m,X_\ast),X_\ast\bigr)
 =\Theta_K\!\left(m\log_2\left(1+\frac{2^p}{m}\right)\right)
 \qquad(p\geq6,\ m\geq2).
\]
The same order holds for fixed pointwise error below one half. Moreover,
$\mathcal A_K$ is a proper subclass of $\mathcal A_K^{\mathrm{rd}}$ on
$X_\ast$ for every permitted precision and every positive block length.
\end{corollary}

\begin{proof}
The previous rounded states remain integers. Thus a common addition still
preserves integer gaps before rounding. Lemma~\ref{lem:rounding-fringe}
extends the lowest uniform-lattice region without introducing a new
boundary. Its error on the added fringe is at most $1/4$, within the old
$1/2$ bound. All local charge inequalities and their telescoping proofs
therefore remain valid. The same codes give the upper bound, and the old
threshold family already satisfies the stronger promise, giving both
lower bounds.

For strict inclusion use the representable operand $-33/4$ and pad with
zeros. At the smallest query $L+8$, its exact sum is $L-1/4$, while its
rounded output is $L$. At every query its output is $x-8$, inside $C_K$.
The operand has six significant bits and normalized exponent three, so
it is available under the stated assumptions. The exact-sum promise fails.
\end{proof}

\subsection{Two thresholds for uniform compositions}

Write $\RN_h$ for nearest-even rounding on the uniform lattice $h\mathbb Z$.
The next lemma concerns uniform rounders on the entire real line;
their application to floating-point arithmetic will be piecewise.

\begin{samepage}
\begin{lemma}[Two-threshold normal form]\label{lem:two-threshold}
Let $d>0$ and all $h_j>0$ be powers of two. Starting from $x_0+di$,
$i\in\mathbb Z$, apply any finite sequence of translations followed by
$\RN_{h_j}$. Denote the resulting function by $G(i)$ and set
\[
 H=\max\{d/2,h_1,\ldots,h_t\},\qquad P=2H/d.
\]
Then $P$ is a positive integer. There exist $c=G(0)$ and integers
$1\leq t_1\leq t_2\leq P$ such that, writing $i=qP+r$ with $0\leq r<P$,
\begin{equation}\label{eq:two-threshold}
 G(i)=c+2Hq+H\bigl(\chi_{r\geq t_1}+\chi_{r\geq t_2}\bigr),
\end{equation}
where $\chi$ denotes an indicator. The thresholds may coincide or equal $P$.
\end{lemma}
\end{samepage}

\begin{proof}
Every uniform rounder commutes with translation by $2H$: this displacement
is a multiple of twice its spacing and preserves tie parity. Hence
$G(i+P)=G(i)+2H$. The function is nondecreasing.

If an operation has spacing $H$, immediately after its first occurrence
all outputs lie in $H\mathbb Z$. A later spacing $H$ has the same property.
For a smaller spacing $h\leq H/2$, translation by any multiple of $H$
preserves rounding, since $2h$ divides $H$. Such an operation preserves
the property that all output differences are multiples of $H$, although
it can change their common residue. The final outputs therefore lie in
the coset $c+H\mathbb Z$. If no operation has spacing $H$, then $H=d/2$;
every step preserves the initial spacing $d$, so $G(i)=c+di$ and the
claim follows with $P=t_1=t_2=1$.

For $0\leq i\leq P$, monotonicity and equivariance give
$c\leq G(i)\leq c+2H$. The coset property leaves only the levels
$c,c+H,c+2H$. Define $t_j=\min\{i\geq1:G(i)\geq c+jH\}$ for
$j=1,2$, allowing the middle level to be skipped.
Equivariance then proves the formula for every integer index, including
negative ones.
\end{proof}

For a restriction to $0\leq i<n$, put $T=\min(P,n)$. Each threshold can
be found by binary search among $0,\ldots,T-1$, with sentinel $T$ if its
level is not reached. If $T<P$, all queries have $q=0$ and $r<T$; clipping
the thresholds at $T$ is therefore harmless. No evaluation outside the
permitted interval is needed. Rebasing a subinterval at its first index
gives the same statement with a new anchor.

\clearpage
\begin{samepage}
\subsection{A constructor polynomial in the number of binades}

\begin{theorem}[Piecewise staircase evaluator]\label{thm:staircase}
Let $n=|X_\ast|$ and let a block of length $m\geq1$ satisfy the rounded-state
$C_K$ promise on $X_\ast$. Its transfer has an exact representation by
at most
\[
 q_m=\min\{n,1+(K-1)m\}
\]
contiguous index intervals, each carrying one two-threshold record.
The representation can be built from the operands without enumerating
$X_\ast$. A query uses $O(\log(q_m+1))$ comparisons and a constant number
of abstract arithmetic operations. A sorted implementation constructs it
using
\begin{equation}\label{eq:staircase-construction}
 O\!\left(Km^2\log(n+1)+Km\log(n+1)\log(Km+2)\right)
\end{equation}
abstract arithmetic, rounding, and comparison operations. Its records
admit a fixed-width description of $O(q_m(p+\log(K+1)))$ bits.
\end{theorem}
\end{samepage}

\begin{proof}
Initially the entire input grid is one affine piece, with $d=4$, $H=2$
and $P=1$. Maintain a partition on which each previous operation has
used a single uniform lattice; the sequence of lattices may differ
between pieces. Each piece consequently has the normal form of
Lemma~\ref{lem:two-threshold} for the corresponding uniform composition.

For an operand $a$, the exact values $F(i)+a$ are nondecreasing. Each
interior binade boundary has one lower-bound cut index. Merge these at
most $K-1$ indices with the existing piece endpoints. On each resulting
subinterval the current operation uses a fixed spacing $h$. At a
boundary itself either adjacent lattice gives the same rounded value;
the lower fringe is covered by Lemma~\ref{lem:rounding-fringe}.

For a new subinterval, retain its parent's maximum spacing and replace
it by $H'=\max(H,h)$. Evaluate its new first output and find the two
thresholds by the bounded binary searches just described. Every probe
uses the old record within its own valid interval, followed by one
uniform rounding. The underlying uniform composition supplies the
normal-form proof even when its extension outside that interval does
not model actual floating-point rounding. Each step adds at most $K-1$
pieces, proving the count.

Two rounded extreme outputs certify the corridor promise by monotonicity.
If either fails, the update is rejected before replacing the old records.
Locating a boundary costs $O(\log(n+1)\log(q+1))$ using the current
piece search. Updating a known subinterval costs $O(\log(n+1))$ probes.
Sorting and merging the cuts is absorbed by the displayed bound. Summing
over all prefixes, with $q\leq1+(K-1)m$, proves
\eqref{eq:staircase-construction}.

For the bit bound, encode piece endpoints and clipped thresholds in
$O(\log(n+1))=O(p)$ bits each. An anchor is one of $KR/2$ floating-point
states. The maximum spacing is a power of two with $O(K)$ possible
exponents and requires $O(\log(K+1))$ bits. Store its exponent, not an
uncompressed binary integer $2^e$. Pad to the public maximum record
count. These fields determine all queries and give the claimed size.
\end{proof}

This code is not asserted to be information optimal: transmitting the
effective operands can use fewer bits. Its purpose is exact evaluation
from a constructed index, with a record bound polynomial in $K$. The
previous residue-class constructor can initialize up to $2^{K-2}$ classes.
Neither abstract operation count measures integer bit costs or GPU time.
The theorem constructs from the original operands and does not provide
an algorithm for merging two descriptions without them.

\section{An evaluator on the entire finite format}
\label{sec:full-format}

The staircase representation also extends beyond positive corridors.
This is an evaluation theorem, not a matching information lower bound.
For this section fix a binary format with precision $p\geq2$, normal
exponents $e_{\min},\ldots,e_{\max}$, and gradual underflow. Put
$E=e_{\max}-e_{\min}+1\geq1$ and
$\delta=2^{e_{\min}-p+1}$. Identify the two signed zeros numerically.
Inputs and operands are finite representable values. Outputs may be
$-\infty$ or $+\infty$; each is absorbing under subsequent finite additions.
NaN inputs, nonfinite operands, exception flags, and zero-sign recovery
are outside this contract. Every addition is correctly rounded separately
to nearest, with ties to even.

\begin{samepage}
\begin{theorem}[Whole-format exact index]\label{thm:full-format}
There are
\[
 N=2^p(E+1)-1
\]
finite numerical input values. For every block of $m\geq0$ finite operands,
its sequential transfer on all these inputs admits an exact index with
at most
\begin{equation}\label{eq:full-piece-bound}
 q_m=\min\{N,2E+1+2Em\}
\end{equation}
contiguous input-rank pieces. Each finite-output piece uses a two-threshold
record; infinite-output pieces are constant. The index can be constructed
without an initial table of all $N$ values. Queries use
$O(\log(E(m+1)+2))$ abstract operations, including input-rank conversion.
A sorted implementation has construction cost
\begin{equation}\label{eq:full-build-bound}
 O\!\left(E+Em^2\log(N+2)
       +Em\log(N+2)\log(E(m+1)+2)\right).
\end{equation}
There is a fixed-width encoding using
$O(q_m(p+\log(E+1)))$ bits. These are arithmetic-operation and information
bounds, not unit-cost assertions about arbitrary-size integers.
\end{theorem}
\end{samepage}

\begin{proof}
Scale all values by $1/\delta$. The subnormal step becomes one. Each
normal binade has $2^{p-1}$ values and spacing $2^j$ for some
$0\leq j<E$. The positive subnormals contribute $2^{p-1}-1$ more values.
Including negatives and one numerical zero gives the displayed $N$.
In increasing order, the finite inputs are the union of $E$ negative
binade progressions, one central progression of subnormals and zero,
and $E$ positive binade progressions. Thus the identity starts with
$2E+1$ affine pieces. Each contains fewer than $2^p$ inputs.

Let $h_{\max}=2^{E-1}$ in the scaled units. The largest finite number
is $M=(2^p-1)h_{\max}$, and the overflow threshold is
\[
 \tau=M+h_{\max}/2.
\]
Nearest-even rounding returns $+\infty$ at $z\geq\tau$ and $-\infty$
at $z\leq-\tau$, including the midpoint ties. Between these tails,
rounding agrees with uniform dyadic rounders on $2E-1$ consecutive
regions. The central region includes both subnormals and the lowest
normal binades. Its spacing is one. The other region boundaries are
\[
 \pm2^{p-1+j},\qquad 1\leq j<E.
\]
Adjacent uniform descriptions agree at each boundary itself.

Suppose the current transfer is indexed. Addition of a finite operand
preserves monotonicity, also on the extended ordered set containing the
two infinities. Partition its exact outputs at the $2(E-1)$ interior
boundaries and the two overflow thresholds. Each has one preimage cut
in the entire ordered input domain. For the negative overflow threshold
use the first exact output strictly greater than $-\tau$; for the
positive threshold use the first output at least $\tau$. Interior cuts
may use the latter convention. These choices include both overflow ties
in their absorbing tails. Existing infinite outputs remain infinite.

After merging these at most $2E$ cuts with old piece endpoints, each new
piece either has one constant infinite output or uses one uniform
rounding lattice. In the latter case its earlier finite history and this
operation form a uniform composition on its original dyadic input
progression. Lemma~\ref{lem:two-threshold} applies with that progression's
spacing $d$. Rebase at the new piece start, update $H$, and recover its
anchor and two clipped thresholds by binary search within the piece.
No history can return from infinity to a finite state, because all
operands are finite. Induction gives \eqref{eq:full-piece-bound}.

Input-rank conversion locates one of the initial progressions and then
uses an exact offset and division by its dyadic spacing. Both this search
and the record search take logarithmically many comparisons. A record
uses \eqref{eq:two-threshold} or returns its constant infinity.
For construction, each region cut requires $O(\log(N+2))$ index probes,
each costing $O(\log(q+1))$. A known new piece requires only
$O(\log(N+2))$ constant-cost probes to its parent record. Sorting and
merging are absorbed by these bounds. Summing over all steps gives
\eqref{eq:full-build-bound}, with $O(E)$ initialization.

Finally, piece endpoints are ranks in a set of size $N$ and anchors
are finite format values or infinities. Each uses
$O(p+\log(E+1))$ bits. The original input spacing and the maximum
rounding spacing have $O(E)$ possible dyadic exponents in the scaled
format; the initial fictitious half-spacing adds only one possibility.
Each clipped threshold is at most its piece length, which is less than
$2^p$ since the initial progression boundaries are retained. Encoding
these fields, the piece count, and padding to $q_m$ proves the bit bound.
\end{proof}

The finite-domain representation can reach $N$ pieces; the theorem does
not guarantee a small index for every length and format. Its bound is
polynomial in the number of binades and the block length, without
requiring enumeration of all significands. The lower bounds on the
corridor families remain consistent with this larger index. Transmitting
the operands is still a smaller general description in many regimes.
The index trades preprocessing and space for queries that do not
replay the block. An efficient merge of two already constructed indices,
and any benefit on a concrete parallel workload, remain separate questions.



\section{Order information with a public numerical multiset}
\label{sec:order}

The threshold family varies operand magnitudes. We now fix the entire
multiset and isolate the information carried by operand order.

\begin{theorem}[Factorial attainability]
\label{thm:factorial}
Let $k\geq 1$ and $R=2^p\geq 64k^2$. There is a public multiset of
$4k$ distinct finite representable operands with permutations realizing
at least $k!$ different transfers on one common set of $k^2$ incoming
markers. Every constructed block has exact real sum zero and satisfies
$F(0)=0$. On the markers, all its exact sums and rounded states lie in
$[3R/4,2R)$. Thus an exact fixed-width description requires at least
$\lceil\log_2(k!)\rceil$ bits even when the multiset is public.
\end{theorem}

\begin{proof}
Index $i,j$ from zero to $k-1$. Set $A=3R/4$ and define
\begin{equation}
\label{eq:permutation-parameters}
\begin{aligned}
 X&=\{A+4\ell:1\leq\ell\leq k^2\},
 &s_i&=R/16+4i,\\
 t_j&=13R/16+4k(j+1),
 &r_j&=R+c_j-t_j,
\end{aligned}
\end{equation}
where $c_j$ is one for even $j$ and four for odd $j$. Given a
permutation $\pi$ of $\{0,\ldots,k-1\}$, concatenate the $k$ slots
\begin{equation}
\label{eq:permutation-slot}
 (+s_{\pi(j)},+r_j,-r_j,-s_{\pi(j)}),\qquad 0\leq j<k.
\end{equation}

All operands are integers of magnitude less than $R$, and hence are
representable. The shifts are distinct and satisfy
\[
 R/16\leq s_i\leq R/8-4.
\]
The radii satisfy $R/8+1\leq r_j\leq 3R/16$ and strictly decrease:
consecutive $t_j$ differ by $4k$, whereas consecutive $c_j$ differ in
magnitude by only three. The two families of positive magnitudes are
disjoint and contain no repetitions. Including their negatives gives
$4k$ distinct operands, each appearing exactly once in every block.
Each slot has exact real sum zero. Starting from zero, its integer
partial sums have magnitude below $R$, so it also returns exactly to
zero in the floating-point format.

Slot $j$ has effective threshold
\begin{equation}
\label{eq:effective-threshold}
 \tau_j=t_j-s_{\pi(j)}=A+4\bigl(k(j+1)-\pi(j)\bigr).
\end{equation}
This threshold occupies one of the $k$ markers numbered $kj+1$ through
$k(j+1)$. The bands are disjoint and strictly ordered. The thresholds
therefore increase with $j$, and their positions determine $\pi$.

For a marker $x$, let $d$ be the parity of the number of thresholds
already crossed. Suppose its state before the slot is $x-d$, as is
true initially. Put $y=x+s_{\pi(j)}$. The range assumptions give
\[
 A+4\leq x\leq 13R/16,
 \qquad 13R/16+4\leq y\leq 15R/16-4.
\]
The first shift is exact and gives $y-d$. Both $y$ and $t_j$ are
divisible by four.

If $x<\tau_j$, then $y\leq t_j-4$. The first sum of the inner pair
is a positive integer at most $R$. Both inner operations are exact
and preserve $y-d$. If $x\geq\tau_j$, then all $j$ earlier thresholds
have been crossed, so $d=j\bmod 2$. For even $j$, the first inner
sum is $R+1+y-t_j$ and rounds down by one. For odd $j$, it is
$R+3+y-t_j$ and rounds up by one. Exact subtraction of $r_j$ gives
$y-1$ or $y$, respectively. The final subtraction of $s_{\pi(j)}$
is also exact, yielding $x-1$ or $x$. The largest exact inner sum is
at most $9R/8-4$, and rounding adds at most one; every exact and
rounded intermediate value remains in $[A,2R)$.

Thus
\begin{equation}
\label{eq:permutation-transfer}
 F_\pi(x)=x-\bigl(\#\{j:\tau_j\leq x\}\bmod 2\bigr).
\end{equation}
As in Theorem~\ref{thm:binomial}, the first threshold in a symmetric
difference distinguishes two transfers. Its band position recovers
the assigned shift. Hence all $k!$ permutations in the family give
different transfer functions.
\end{proof}

For $m=4k$ and $R\geq 4m^2$, this gives $\Theta(m\log m)$ bits of
additional order information as $k$ grows. The matching upper bound
is an index among at most $m!$ orders of the public multiset. This is
a worst-case result over numerical multisets in the stated precision
regime, not a lower bound for every multiset. Common zero padding
preserves the $k!$ witnesses at longer lengths, but does not permit
replacing $k$ by an arbitrarily larger $m$ in the asymptotic bound.
For $m=256$ and $k=64$, the exact deterministic lower bound is $296$
bits whenever the precision condition holds. This conditional-multiset
bound concerns a different interface from the unrestricted one in
Section~\ref{sec:threshold}.

\section{Pointwise bounded-error summaries}
\label{sec:random}

\subsection{Entropy lemma}

Write
\[
 h_2(\varepsilon)=-\varepsilon\log_2\varepsilon
 -(1-\varepsilon)\log_2(1-\varepsilon),
\]
with $0\log_2 0=0$. The following standard entropy argument separates
the information inequality from the arithmetic reductions.

\begin{lemma}[Conditional bit recovery]
\label{lem:entropy}
Let $B_1,\ldots,B_q$ be independent uniform bits, let $U$ be public
randomness independent of them, and let $M$ be a $b$-bit message.
Suppose that, for every $t$, an estimator given
$M,U,B_1,\ldots,B_{t-1}$ recovers $B_t$ with error probability at
most $\varepsilon$, where $0\leq\varepsilon<1/2$. Then
\begin{equation}
\label{eq:entropy-bound}
 b\geq\bigl(1-h_2(\varepsilon)\bigr)q.
\end{equation}
\end{lemma}

\begin{proof}
Let $E_t$ indicate whether the estimator is wrong. Given the
estimator's inputs, $E_t$ determines $B_t$. The binary entropy bound
therefore gives
\[
 H(B_t\mid M,U,B_1,\ldots,B_{t-1})\leq h_2(\varepsilon).
\]
The chain rule yields
\[
 H(B_1,\ldots,B_q\mid M,U)\leq qh_2(\varepsilon).
\]
Independence from $U$ gives $H(B_1,\ldots,B_q\mid U)=q$. Consequently,
\[
 I(B_1,\ldots,B_q;M\mid U)
 \geq q\bigl(1-h_2(\varepsilon)\bigr).
\]
The same mutual information is at most $H(M\mid U)\leq b$.
\end{proof}

In the applications below, the query recovering a bit depends only on
earlier bits and public parameters, not on $M$ or $U$. The pointwise
guarantee for each fixed block-query pair thus bounds the averaged
decoding error in Lemma~\ref{lem:entropy}. No simultaneous guarantee
over all queries of one random summary is needed.

\subsection{Threshold choices}
\label{subsec:random-thresholds}

Let $N$ denote the size of either threshold grid, and choose
\[
 1\leq k\leq\min\{\lfloor m/2\rfloor,\lfloor N/2\rfloor\},
 \qquad q=\lfloor\log_2(N/k)\rfloor\geq 1.
\]
Partition the first $k2^q$ markers into $k$ consecutive groups of
$2^q$ markers. In each group choose one threshold whose rank is an
independent uniform $q$-bit number. Theorem~\ref{thm:binomial}
constructs a block for every such choice.

Order the bits first by group and then, within a rank, from most to
least significant. Given the higher bits of a rank, its possible
values form a dyadic interval. Query the marker at the last rank in
the lower half of that interval. After removing the known parity
of earlier groups, the transfer output says whether the selected
threshold is at most this marker, and hence determines the next bit.
Lemma~\ref{lem:entropy} gives
\begin{equation}
\label{eq:random-threshold-bound}
 b\geq\bigl(1-h_2(\varepsilon)\bigr)
       k\lfloor\log_2(N/k)\rfloor.
\end{equation}

For the corridor grid, take
$k=\min\{\lfloor m/2\rfloor,\lfloor N_\ast/2\rfloor\}$.
This yields $\Omega(m\log_2(1+R/m))$ before saturation and
$\Omega(R)$ afterward, completing the bounded-error part of
Theorem~\ref{thm:corridor}. Its constants may depend on the fixed
error $\varepsilon$. Together with the operand-replay upper bound,
the same argument gives $\Theta(mp)$ in the regime of
Corollary~\ref{cor:short}.

For $m\geq 2N$, a direct construction covers saturation. Encode an
arbitrary vector of $N$ bits by selecting a threshold wherever
consecutive bits differ, taking the bit before the first marker to
be zero. At most $N$ threshold pairs are needed. Common zero padding
fixes the length, and the output at each marker reveals its bit.
Lemma~\ref{lem:entropy} gives
$b\geq(1-h_2(\varepsilon))N$. In the deterministic case, this
subfamily requires exactly $N$ bits.

\subsection{Permutation ranks}
\label{subsec:random-permutations}

In Theorem~\ref{thm:factorial}, order the remaining shifts at slot $j$
by their effective threshold. There are $k-j$ choices. Set
$q_j=\lfloor\log_2(k-j)\rfloor$, retain the first $2^{q_j}$ ranks,
and choose a rank using independent uniform $q_j$ bits. Earlier
choices determine the remaining ordered set, so every chosen rank
has an available shift. Distinct rank sequences produce distinct
permutations.

Given the earlier slots and the higher bits of the current rank,
query the marker at the last remaining threshold in the lower half
of the current dyadic rank interval. It is an allowed marker in
slot $j$'s band. All previous slots have crossed, no later slot has
crossed, and the current output parity determines the next rank bit.
Applying Lemma~\ref{lem:entropy} to the concatenated bits gives
\begin{equation}
\label{eq:random-permutation-bound}
\begin{aligned}
 b&\geq\bigl(1-h_2(\varepsilon)\bigr)
          \sum_{u=1}^k\lfloor\log_2 u\rfloor\\
  &\geq\bigl(1-h_2(\varepsilon)\bigr)\bigl(\log_2(k!)-k\bigr).
\end{aligned}
\end{equation}
Thus the $\Theta(m\log m)$ order-information bound also holds for
every fixed pointwise error below one half, in the same precision
regime. The second expression in \eqref{eq:random-permutation-bound}
may be uninformative for small $k$; the exact sum is always
nonnegative.

All these bit-recovery arguments also apply when, with probability
at least $1-\varepsilon$, the numerical answer differs from the
correct result by strictly less than $1/2$. The two possible answers
at a query differ by one, so their permitted answer intervals are
disjoint. An absolute error allowance of one destroys these
particular witnesses: returning $x$ then succeeds without any
block-dependent message.

\section{Validation and reproducibility}\label{sec:validation}

\subsection{Baselines and evaluation criteria}

We compare description sizes. We use three exact baselines:
transmitting the $m$ operand words and replaying the block after $x$ arrives;
storing every query result; and encoding a nondecreasing output sequence by a
monotone bitvector. For the constructed families, subset and permutation ranks
give the corresponding information upper bounds. These comparisons concern
all representations obeying the access contract, rather than the properties
of one implementation.

The reference implementation uses exact rational arithmetic to construct
sparse gap maps and fixed-width byte messages. A full-table oracle checks the
same ordered specification on small domains, and the monotone fallback is
implemented. The serialized size is distinguished from Python object memory.
Native CPU and GPU runs test arithmetic correctness, not throughput. The
reference is not claimed to outperform a native sequential fold or to provide
a GPU optimization.

The artifact also retains an exploratory arithmetic implementation of the
Kapre-DeHon iterative scheme~\cite{ref4}. It does not reproduce their FPGA
implementation or benchmark a current fastest kernel. Its correction counts
refer to a different data-access contract and are not an optimality baseline
for the standalone summary.

\subsection{Completed checks}

Table~\ref{tab:validation} lists the completed checks. The saved hardware
outputs are kept separate from subsequent comparisons against those outputs;
the latter are not counted as additional device runs.

\begin{table}[tbp]
\small
\renewcommand{\arraystretch}{1.13}
\caption{Completed validation. Counts refer to executed checks, not planned experiments.}
\label{tab:validation}
\begin{tabularx}{\textwidth}{@{}P{.26\textwidth}P{.45\textwidth}Y@{}}
\toprule
Check & Coverage & Outcome \\
\midrule
Exact threshold enumeration & 16,452 subsets; 229,768 marker folds & Expected cardinalities \\
Exact permutation enumeration & 873 permutations; 29,367 marker folds & Exactly $k!$ signatures \\
Factoradic query reduction & 1,367 codewords; 12,746 bit queries & Every bit recovered \\
Corridor block enumeration & 11,652 admissible blocks; 60,036 steps & Conservative invariant holds \\
Original local symbolic checks & 196 counterexample queries & All UNSAT \\
Independent symbolic checks & 7 local obligations, including the weighted $4m$ bound & All UNSAT \\
Independent finite-format arithmetic & 59,528 additions at $p=2,3,4$, including exceptional numerical cases & Oracle checks pass \\
Independent corridor exploration & 142,340 reachable-state transitions & Support and $4m$ bounds hold \\
Native CPU arithmetic & 120,925 queries in binary16, binary32, and binary64 & No output mismatches \\
Actual Apple M5 Pro GPU arithmetic & 27,713 queries in binary32 & No output mismatches \\
Reference constructor & 144 blocks; 2,280 prefixes; 107,920 full-table comparisons & Pass with $4m$ assertion \\
Independent reference cross-check & 84 blocks; 18,240 queries; 3,444 serialization round trips & Exact agreement \\
\bottomrule
\end{tabularx}
\end{table}

The seven independent symbolic obligations cover monotone rounding,
residue-potential decrease away from the boundary, a crossing's gap change of
at most one, preservation of a gap of four away from a crossing, both weighted
local inequalities, and the threshold-pair transition. They are machine
checks of local lemmas, not a proof-assistant formalization of the complete
theorems. Earlier unsplit symbolic searches that timed out remain recorded as
unsuccessful attempts.

The independent arithmetic oracle includes signed zeros, gradual underflow,
and overflow to test numerical cases beyond the witness trajectories.
Finite enumeration does not establish a theorem at arbitrary precision; the
general arguments are in Sections~\ref{sec:threshold}--\ref{sec:random}.
Likewise, the reachable-state count records explored transitions and does not
claim to exhaust all real operand sequences.

The CPU records include strict rounding settings and scalar floating-add
instructions. The GPU records include the device, compilation options, input
buffers, and actual output bits. They contain 7,560 threshold queries and
20,153 permutation queries, totaling 27,713. The later sparse-constructor
comparison with the 20,153 permutation outputs reuses those saved bits.
Software arithmetic models are recorded separately from these device runs.

The fixed-$K$ extension was checked separately using exact arithmetic. A
uniform-lattice enumeration covered 347,376 cases, and 1,339,284 local
endpoint pairs checked the boundary charge. Two independently implemented
rounders agreed on 8,608 individual exact inputs. A further 600 valid blocks
gave 20,520 checked prefixes and 430,920 full-table query comparisons, with
up to six binades. An independently written rational oracle checked
11,169,226 uniform-lattice cases, 2,509,836 boundary pairs, and 6,144 block
prefixes with 122,880 queries; its finite encoding test made 11,181 round
trips. These counts describe executed finite checks, not a formal proof of
all parameters.

The new sparse evaluator passed 240 block tests, comprising 8,208 prefixes
and 103,968 full-table comparisons, and rejected three corridor violations
without changing its state. Additional tests with nonbinary rational
addends checked the rounding implementation but are not claimed as
floating-point operand blocks. For $K=3$ and $p=24,53$, the representable
operand $3R/2+2$ produces alternating gaps eight and zero on $X_\ast$.
The evaluator represents this map by two spacing-eight classes with no
deviation entries. At each precision, 518 selected queries agreed with
direct exact rounding. The domain sizes are the same as in
Table~\ref{tab:serialized}; these are evaluator tests, not new serialized
messages or hardware measurements.

\subsection{Staircase indices and the complete format}

The rounded-state relaxation was checked on 1,495,456 local endpoint
pairs, including 2,776 lower-fringe pairs, and 560 blocks with 19,152
prefixes and 402,192 full-grid comparisons. Of those prefixes, 562
actually violate the old exact-sum promise while satisfying the new one.
The proof, rather than this finite enumeration, establishes the relaxation
at every permitted precision.

An independently written normal-form checker tested 87,165 short uniform
compositions and 160 longer chains, making 904,342 query comparisons.
Its separate corridor piece updater made 402,192 full-grid comparisons.
The main staircase implementation passed 128 blocks, 4,096 prefixes,
and 339,968 full-grid comparisons; a separately written cross-check added
104,580 comparisons. These tests cover coincident thresholds, clipped
periods, negative index rebasing in the uniform lemma, and preservation
of a coset after a finer rounding.

For the whole-format reference, an enumerated-neighbour oracle checked
30,018 exact sums. All finite inputs were then checked after every prefix
of 984 exhaustively generated short blocks and 400 additional blocks:
10,404 prefixes and 912,216 query comparisons in total. Tests cover
negative binade boundaries, gradual underflow, both overflow midpoint
ties, and the persistence of infinity after a later finite operand.
A wrong lattice-side choice at a negative boundary was detected in the
prototype and corrected before these final runs.

Table~\ref{tab:full-format} gives two additional whole-format indices.
The same eight-operand pattern in each format was
\[
 (\delta,1,-1,2^p,1,-2^p,M,-M),
\]
where $\delta$ is its minimum positive subnormal and $M$ its largest
finite value. These indices cover the entire finite numerical input
domain by Theorem~\ref{thm:full-format}; their construction did not
enumerate that domain. Sampled inputs include piece endpoints and random
input ranks. Their outputs were compared against actual CPU additions
using NumPy 2.3.5 on arm64 macOS, with an explicit binary32 or binary64
dtype at every addition. Subnormal preservation was checked. The input,
operand, and output bits are retained. The 23,056 sampled queries used
184,448 separate CPU additions and produced no numerical mismatches.
These are arithmetic correctness checks, not timings or new GPU runs.

\begin{table}[tbp]
\caption{Whole-format indices for an eight-operand block. Counts are records,
not bytes or speed measurements. Signed zeros count as one input.}
\label{tab:full-format}
\small
\begin{tabularx}{\textwidth}{@{}Yrr@{}}
\toprule
Quantity & binary32 & binary64 \\
\midrule
Normal binades $E$ & 254 & 2,046 \\
Finite numerical inputs $N$ & 4,278,190,079 & 18,437,736,874,454,810,623 \\
Constructed records & 1,256 & 7,992 \\
Sampled native-CPU queries & 6,165 & 16,891 \\
Numerical mismatches & 0 & 0 \\
\bottomrule
\end{tabularx}
\end{table}

The corridor constructor was also tested at $p=24,53$ and $K=3,20,64$.
Each of six examples used 33 records on grids of 1,048,574 or
562,949,953,421,310 queries, with 322 sampled queries checked against
exact arithmetic. The record count is not a comparison against the
smallest possible encoding; some blocks have further compressible outputs.

\subsection{Serialized examples on large domains}

For two length-256 threshold blocks with 128 selected thresholds, the sparse
constructor was run on the corridor grid at precisions 24 and 53. The
representations were constructed without traversing the full query domain.
Table~\ref{tab:serialized} gives their sizes and construction counts.

\begin{table}[tbp]
\caption{Serialized corridor examples with $m=256$. The public budgets use the conservative $5m$ serializer.}
\label{tab:serialized}
\small
\renewcommand{\arraystretch}{1.18}
\begin{tabularx}{\textwidth}{@{}Yrr@{}}
\toprule
Quantity & $p=24$ & $p=53$ \\
\midrule
Incoming states & 1,048,574 & 562,949,953,421,310 \\
Actual payload bits & 2,602 & 6,343 \\
Public worst-case bit budget & 4,645 & 12,099 \\
Serialized bytes & 581 & 1,513 \\
Rounding evaluations & 33,536 & 33,536 \\
\bottomrule
\end{tabularx}
\end{table}

The payload depends on the block, but every message is padded to the public
worst-case budget. The byte counts include final byte padding and exclude
shared format, grid, and length metadata. Both examples use the conservative
$5m$ magnitude budget retained by the serializer; the $4m$ theorem permits a
tighter information bound. At each precision, 642 sampled queries agree with
direct exact folds. Full-table tests on small domains, rather than those
samples alone, check the general implementation path. The table reports sizes
and operation counts, not running time or peak memory.

\section{Limitations and open questions}\label{sec:limitations}

\subsection{The corridor and rounding hypotheses}

The original spacing-four invariant fails outside two binades. At $p=7$,
let $R=128$ and apply the exactly representable operand 512 to the six markers
\[
 72,76,80,84,88,92.
\]
The rounded outputs are
\[
 584,592,592,592,600,608.
\]
Their gap deviations from four are $4,-4,-4,4,4$, with absolute sum 20 and
support five. Both exceed the one-step bounds in
Theorem~\ref{thm:variation}. The failure already occurs in three binades:
at $p=8$, the representable operand 386 maps $X_\ast=\{136,140,\ldots,188\}$
to gaps alternating eight and zero. All 13 gaps deviate from four and their
total deviation is 52 after one addition. Partitioning $X_\ast$ into its
two residue classes modulo eight removes this apparent obstruction to
Theorem~\ref{thm:fixed-K}; it does not restore the original spacing-four
support estimate.

The boundary count in Lemma~\ref{lem:crossing-Q} is also necessary for that
particular charge bound. At $p=6$, $K=4$, $Q=16$, exact endpoints 51 and
291 round to 51 and 288. The gap changes from 240 to 237, giving charge
$3+15-0=18>Q$. The endpoints strictly cross three boundaries, so the
proved bound $Qc$ still holds.

The tie rule also matters. At $p=4$, the block $(+9,-9)$ maps $x=8$ to seven
under ties-to-even and to nine under ties away from zero. The threshold proof
cannot simply substitute another rounding mode. Moreover, $(+9,-9)$ and
$(+5,-5)$ both have real sum zero and return zero from zero, but map eight to
seven and eight, respectively, under ties-to-even. Neither the correctly
rounded real sum nor the zero-start fold determines the required transfer.

The exact-sum condition is strictly stronger than the rounded-state
condition. At $p=7$, the operand $-33/4$ sends the smallest query 72 to
the exact sum 63.75 and the rounded state 64. It violates the former
promise but satisfies the latter on the entire query grid.
Corollary~\ref{cor:rounded-corridor} proves that the same information and
variation bounds still apply. The original sparse implementation retains
its stricter admission check; the new staircase constructor accepts the
relaxed class.

The zero-start identity in Theorem~\ref{thm:factorial} applies to its
constructed orders, not every permutation of the public multiset. For
$p=11$ and $k=5$, list all positive operands as
\[
 128,132,136,140,144,365,348,325,308,285
\]
and then append their negatives in the same order. The block returns one
from zero. The slot-form orders used in the proof return zero.

\subsection{The access contract and composition}

Interaction changes the information problem. A decoder can send $x$ to the
block owner and receive $F_a(x)$ using two floating-point words. If $x$ is
already known during encoding, one output word suffices. These alternatives
lie outside the unknown-entry one-way contract. The lower bounds also permit
algorithms that revisit the original operands, use distributed storage charged
elsewhere, or perform iterative speculative correction.

The corridor sparse representation is an evaluator on its specified grid. That grid
need not contain the outputs of another block, so an exact associative scan
format would require an additional closure argument. The implementation also
checks a narrower starting domain than the abstract spacing-four lemma; the
grid $X_\ast$ satisfies both. Its construction bound counts abstract arithmetic
operations and does not establish a bit-complexity or GPU-work bound.

\subsection{Unresolved questions}

The matching fixed-$K$ information result leaves open a tight joint
dependence on length, precision, and the number of binades when $K$ grows.
The new staircase index has polynomial record dependence, but is generally
larger than the operand list and does not close that information gap.
The whole-format theorem removes positivity for evaluation; it identifies
signed zeros numerically and omits NaNs, nonfinite operands, and exception
flags. An efficient algorithm for merging two constructed indices remains
open. A separate systems question is whether constructing
and applying summaries pays for its cost on a concrete workload. Answering it
requires optimized implementations, a specified memory-access contract, and
comparable timing measurements. Our results do not settle either question.

\section{Conclusion}\label{sec:conclusion}

Ordered rounded additions realize many distinct functions of the incoming
accumulator, even when the real sums and zero-start folds agree. Threshold
and permutation constructions turn this distinction into description lower
bounds. On a common two-binade query domain, a residue potential bounds total
gap variation by $4m$, yielding a sparse evaluator and a matching
$\Theta\bigl(m\log_2(1+2^p/m)\bigr)$ bound. A greatest-common-divisor
potential and a partition of the same query grid extend that order to every
fixed number of adjacent positive binades, and the same bounds hold under
the weaker rounded-state promise. A different two-threshold normal form
gives an exact index with $O(mK)$ corridor pieces and $O(E(m+1))$ pieces
over the entire finite format, including negative values, subnormals,
and absorbing infinities. These evaluation indices need not minimize
description size. Standard entropy arguments extend the lower bounds to
pointwise bounded error. None of these results is a running-time lower
bound for unrestricted parallel algorithms or a measured speedup.

\appendix
\section{Reproduction}\label{sec:artifact}

The accompanying research artifact, version~1.0.0, is available at
\url{https://doi.org/10.5281/zenodo.24044237}. Its \path{README.txt}
describes the package, \path{REPRODUCE.txt} gives the dependencies and
commands, and \path{EVIDENCE_MAP.json} maps the paper's validation results
to code and data. The source code is in \path{research/experiments/},
\path{research/adversarial/}, and \path{research/extensions/}; recorded
outputs and environment information are in \path{research/results/} and
the relevant extension directories.

The following commands run from \path{research/}. The symbolic checks and
the independent reference cross-check require the Z3 Python package. The
boundary-attack, corridor-summary, and retained-record verification commands
use the standard library.
\begin{quote}\footnotesize\ttfamily
python3 -B adversarial/root/boundary\_attacks.py\\
python3 -B adversarial/math/independent\_audit.py\\
python3 -B adversarial/math/reference\_crosscheck.py\\
python3 -B experiments/check\_corridor\_summary.py\\
python3 -B experiments/verify\_artifact.py --public-bundle
\end{quote}
The last command checks retained code and data without rerunning the GPU
experiments. Device build and execution instructions are in
\path{REPRODUCE.txt} at the package root. Saved outputs support checks on
machines without the original GPU but do not replace an independent device
run. The independent reference cross-check used seed 89623471; the other
seeds are recorded in their scripts and result files. The fixed-$K$
extension and its separate checks are in
\path{research/extensions/multibinade_2026_11_07/}. Its two main commands
are \texttt{python3 -B check\_multibinade.py} and
\texttt{python3 -B check\_sparse.py}, run from that directory. The independent
oracle is \path{independent_review/check_independent.py}.

The rounded-state, staircase, and full-format extensions are in
\path{research/extensions/strengthening_2026_11_07/}. Their principal
commands, run from that directory, are
\begin{quote}\small\ttfamily
python3 -B check\_staircase.py\\
python3 -B independent\_math/check\_rounded\_corridor.py\\
python3 -B independent\_math/check\_uniform\_pieces.py\\
python3 -B independent\_math/check\_full\_format.py\\
python3 -B check\_full\_native.py
\end{quote}
The last command requires NumPy and executes new CPU arithmetic; it does
not use a GPU.

The historical comparison in Section~\ref{sec:related} used the complete
Reiser-Knuth original, its corrections, the corrected 2010 chapter and
addendum, and Diamond's article. Knuth's exercise~4.2.2--22 and its answer
were compared between the second edition (pp.~229 and~573) and the third
edition (pp.~244 and~616). Both give the fixed-multiplier stabilization
result; the later version explicitly excludes $v=0$ and revises a
boundary-case argument. This comparison does not certify an exhaustive
priority search.

\begin{samepage}
\section*{Use of AI tools}

Generative AI assisted with developing and checking mathematical arguments,
writing and testing code, comparing the literature, and editing the
manuscript; the author takes responsibility for its contents.
\end{samepage}

\bibliographystyle{amsplain}
\providecommand{\bysame}{\leavevmode\hbox to3em{\hrulefill}\thinspace}
\providecommand{\MR}{\relax\ifhmode\unskip\space\fi MR }
\providecommand{\MRhref}[2]{%
  \href{http://www.ams.org/mathscinet-getitem?mr=#1}{#2}
}
\providecommand{\href}[2]{#2}
\begin{thebibliography}{10}

\bibitem{ref28}
Roberto Bagnara, Abramo Bagnara, Fabio Biselli, Michele Chiari, and Roberta
  Gori, \emph{Correct approximation of {IEEE 754} floating-point arithmetic for
  program verification}, Constraints \textbf{27} (2022), 29--69,
  \url{https://doi.org/10.1007/s10601-021-09322-9}.

\bibitem{ref3}
Julien Brandoit, Arthur Fyon, Thomas Braipson, Tom Clara, Florent {De Geeter},
  Pierre Sacr{\'e}, Damien Ernst, and Guillaume Drion, \emph{Length-independent
  state tracking under a parallel scan}, arXiv:2609.32447v1, 2026,
  \url{https://doi.org/10.48550/arXiv.2609.32447}.

\bibitem{ref10}
S{\'i}lvia Casacuberta, Michael Shoemate, Salil Vadhan, and Connor Wagaman,
  \emph{Widespread underestimation of sensitivity in differentially private
  libraries and how to fix it}, Proceedings of the 2022 {ACM SIGSAC} Conference
  on Computer and Communications Security ({CCS} '22), ACM, 2022, Full version:
  arXiv:2207.10635v2; \url{https://doi.org/10.1145/3548606.3560708}.

\bibitem{ref16}
Amit Chakrabarti, Graham Cormode, Ranganath Kondapally, and Andrew McGregor,
  \emph{Information cost tradeoffs for augmented index and streaming language
  recognition}, arXiv:1004.3304v1, 2010,
  \url{https://doi.org/10.48550/arXiv.1004.3304}.

\bibitem{ref19}
Harold~G. Diamond, \emph{Stability of rounded off inverses under iteration},
  Mathematics of Computation \textbf{32} (1978), no.~141, 227--232,
  \url{https://doi.org/10.1090/S0025-5718-1978-0461879-7}.

\bibitem{ref24}
Roger~B. Eggleton, \emph{The well-rounded linear function}, The Electronic
  Journal of Combinatorics \textbf{8} (2001), no.~2, R6,
  \url{https://doi.org/10.37236/1605}.

\bibitem{ref7}
\emph{Errata}, Information Processing Letters \textbf{3}, no.~5, p.~164, 1975,
  First correction block: {Reiser and Knuth} (1975);
  \url{https://doi.org/10.1016/0020-0190(75)90031-9}.

\bibitem{ref14}
Luca Foschini, John Hershberger, and Subhash Suri, \emph{On the complexity of
  time-dependent shortest paths}, Proceedings of the Twenty-Second Annual
  {ACM-SIAM} Symposium on Discrete Algorithms ({SODA} 2011), SIAM, 2011,
  \url{https://doi.org/10.1137/1.9781611973082.27}, pp.~327--341.

\bibitem{ref21}
Robert~Joseph George, Will Adkisson, and Anima Anandkumar, \emph{{FloatLib}:
  Verified floating-point arithmetic in {Lean}}, 2026, arXiv:2609.19352v1
  [cs.LO], September 16, 2026; \url{https://arxiv.org/abs/2609.19352v1}.

\bibitem{ref8}
David Goldberg, \emph{What every computer scientist should know about
  floating-point arithmetic}, ACM Computing Surveys \textbf{23} (1991), no.~1,
  5--48, {Oracle} reprint;
  \url{https://docs.oracle.com/cd/E19957-01/806-3568/ncg_goldberg.html}.

\bibitem{ref26}
Monika Henzinger, Stefan Neumann, Harald R{\"a}cke, and Stefan Schmid,
  \emph{Dynamic maintenance of monotone dynamic programs and applications},
  40th International Symposium on Theoretical Aspects of Computer Science
  (STACS 2023), Leibniz International Proceedings in Informatics (LIPIcs), vol.
  254, Schloss Dagstuhl -- Leibniz-Zentrum f{\"u}r Informatik, 2023,
  \url{https://doi.org/10.4230/LIPIcs.STACS.2023.36}, pp.~36:1--36:16.

\bibitem{ref4}
Nachiket Kapre and Andr{\'e} DeHon, \emph{Optimistic parallelization of
  floating-point accumulation}, 18th {IEEE} Symposium on Computer Arithmetic
  ({ARITH} 2007), IEEE, 2007,
  \url{https://nachiket.github.io/publications/fp-accum_arith-2007.pdf}.

\bibitem{ref23}
Donald~E. Knuth, \emph{The art of computer programming, volume 2: Seminumerical
  algorithms}, 2nd ed., Addison-Wesley, Reading, MA, 1981, Exercise 4.2.2--22,
  p.~229; answer, p.~573;
  \url{https://www-cs-faculty.stanford.edu/~knuth/taocp.html}.

\bibitem{ref20}
\bysame, \emph{The art of computer programming, volume 2: Seminumerical
  algorithms}, 3rd ed., Addison-Wesley, Reading, MA, 1997, Exercises 4.2.2--22
  and 4.2.2--28, pp.~244--245; solutions, pp.~616--617;
  \url{https://www-cs-faculty.stanford.edu/~knuth/taocp.html}.

\bibitem{ref1}
Richard~E. Ladner and Michael~J. Fischer, \emph{Parallel prefix computation},
  Journal of the ACM \textbf{27} (1980), no.~4, 831--838,
  \url{https://doi.org/10.1145/322217.322232}.

\bibitem{ref25}
Jeffrey~C. Lagarias, Takumi Murayama, and D.~Harry Richman, \emph{Dilated floor
  functions that commute}, The American Mathematical Monthly \textbf{123}
  (2016), no.~10, 1033--1038,
  \url{https://doi.org/10.4169/amer.math.monthly.123.10.1033}.

\bibitem{ref27}
Wonyeol Lee, Rahul Sharma, and Alex Aiken, \emph{Verifying bit-manipulations of
  floating-point}, Proceedings of the 37th ACM SIGPLAN Conference on
  Programming Language Design and Implementation, PLDI '16, Association for
  Computing Machinery, 2016, \url{https://doi.org/10.1145/2908080.2908107},
  pp.~70--84.

\bibitem{ref17}
Laura Monroe and Vanessa Job, \emph{Computationally inequivalent summations and
  their parenthetic forms}, arXiv:2005.05387v1, 2020,
  \url{https://doi.org/10.48550/arXiv.2005.05387}.

\bibitem{ref2}
Todd Mytkowicz, Madanlal Musuvathi, and Wolfram Schulte, \emph{Data-parallel
  finite-state machines}, Proceedings of the Nineteenth International
  Conference on Architectural Support for Programming Languages and Operating
  Systems ({ASPLOS} '14), ACM, 2014,
  \url{https://doi.org/10.1145/2541940.2541988}.

\bibitem{ref12}
Franz Nowak, Ryan Cotterell, and Reda Boumasmoud, \emph{An algebraic view of
  the expressivity of recurrent language models}, arXiv:2606.01765v3, 2026,
  Accepted at the 43rd International Conference on Machine Learning ({ICML}
  2026); \url{https://doi.org/10.48550/arXiv.2606.01765}.

\bibitem{ref13}
\bysame, \emph{A compositional theory of causally masked transformers},
  arXiv:2607.26988v1, 2026, \url{https://doi.org/10.48550/arXiv.2607.26988}.

\bibitem{ref5}
Karl Papadantonakis, Nachiket Kapre, Stephanie Chan, and Andr{\'e} DeHon,
  \emph{Pipelining saturated accumulation}, IEEE Transactions on Computers
  \textbf{58} (2009), no.~2, 208--219,
  \url{https://doi.org/10.1109/TC.2008.110}.

\bibitem{ref6}
John~F. Reiser and Donald~E. Knuth, \emph{Evading the drift in floating-point
  addition}, Information Processing Letters \textbf{3} (1975), no.~3, 84--87,
  \url{https://doi.org/10.1016/0020-0190(75)90022-8}.

\bibitem{ref18}
\bysame, \emph{Evading the drift in floating-point addition}, Selected Papers
  on Design of Algorithms, CSLI Lecture Notes, vol. 191, Center for the Study
  of Language and Information, Stanford, CA, 2010, Corrected reprint with
  addendum on p.~348; \url{https://www-cs-faculty.stanford.edu/~knuth/da.html},
  pp.~341--348.

\bibitem{ref22}
Brett Saiki, Bill Zorn, Cynthia Richey, and Zachary Tatlock, \emph{When double
  rounding is correct}, 2026, arXiv:2610.09005v2 [cs.PL], October 15, 2026;
  \url{https://arxiv.org/abs/2610.09005v2}.

\bibitem{ref15}
Sebastiano Vigna, \emph{Quasi-succinct indices}, Proceedings of the 6th {ACM}
  International Conference on Web Search and Data Mining ({WSDM} '13), ACM,
  2013, \url{https://vigna.di.unimi.it/ftp/papers/QuasiSuccinctIndices.pdf},
  pp.~83--92.

\bibitem{ref9}
V{\'a}clav Vor{\'a}{\v{c}}ek and Matthias Hein, \emph{Sound randomized
  smoothing in floating-point arithmetics}, International Conference on
  Learning Representations ({ICLR} 2023), 2023, Full version:
  arXiv:2207.07209v2; \url{https://arxiv.org/abs/2207.07209v2}.

\bibitem{ref11}
Peichen Xie, Yanjie Gao, Yang Wang, and Jilong Xue, \emph{Revealing
  floating-point accumulation orders in software/hardware implementations},
  2025 {USENIX} Annual Technical Conference ({USENIX ATC} 25), USENIX
  Association, 2025,
  \url{https://www.usenix.org/conference/atc25/presentation/xie},
  pp.~1425--1440.

\end{thebibliography}


\end{document}
